\documentclass[a4paper,11pt]{article}

\usepackage[utf8]{inputenc}
\usepackage[T1]{fontenc}
\usepackage[ngerman]{babel}
\usepackage{amsmath}
\usepackage{amssymb}
\usepackage{xcolor}
\usepackage{array}
\usepackage{booktabs}
\usepackage{geometry}
\geometry{a4paper,margin=2cm}

\definecolor{accent}{HTML}{1A4E8A}
\definecolor{accent2}{HTML}{8A5A00}
\definecolor{accent3}{HTML}{2E6E3A}
\definecolor{boxbg}{HTML}{F2F6FB}
\definecolor{rulecol}{HTML}{B0B0B0}

\newcommand{\aufg}[2]{%
  \par\vspace{8pt}
  {\large\bfseries\color{accent} Aufgabe #1}\hfill{\small\color{accent2}\bfseries #2}\par
  \noindent\textcolor{accent}{\rule{\linewidth}{0.8pt}}\par\vspace{3pt}}
\newcommand{\teil}[2]{\par\vspace{5pt}{\bfseries\color{accent2}(#1)}~\textit{#2}\par\vspace{1pt}}
\newcommand{\erg}[1]{\fcolorbox{accent3}{boxbg}{$\displaystyle #1$}}

\setlength{\parindent}{0pt}
\setlength{\parskip}{3pt}

\begin{document}

\begin{center}
  {\Large\bfseries\color{accent} Numerische Methoden -- Probeklausur 00}\\[1pt]
  {\normalsize\bfseries Musterlösung mit vollständigem Rechenweg}\\[2pt]
  {\small Prüfung von Prof.\ Dr.-Ing.\ Mark Schutera \;·\; DHBW Ravensburg \;·\; 25 Punkte}
\end{center}
\noindent\textcolor{rulecol}{\rule{\linewidth}{0.4pt}}

{\footnotesize Alle Zwischen- und Endergebnisse wurden nachgerechnet. Endergebnisse sind
\erg{\text{umrandet}}. Brüche bleiben stehen; es wird kein Taschenrechner benötigt.}

% =====================================================================
\aufg{1}{Gleitkomma, Gitter-Suche \& Fehleranalyse \hfill 16 Punkte}

\teil{a}{Gleitkommasystem, Maschinenepsilon, Auslöschung \hfill 4 P}

\textbf{Trade-off Mantisse vs.\ Exponent.} Bei fester Bitzahl gilt:
\begin{itemize}\setlength{\itemsep}{1pt}
  \item Ein zusätzliches \textbf{Mantissenbit} verbessert die \emph{Auflösung}
        (resolution): das Maschinenepsilon halbiert sich, benachbarte Zahlen rücken
        enger zusammen, der relative Rundungsfehler sinkt. Der Wertebereich bleibt gleich.
  \item Ein zusätzliches \textbf{Exponentenbit} vergrößert den \emph{Wertebereich}
        (dynamic range): die größte darstellbare Zahl wächst, die kleinste sinkt
        (mehr Größenordnungen). Die Auflösung pro Größenordnung bleibt gleich.
\end{itemize}
Es ist ein Tausch: feine Auflösung \emph{oder} großer Wertebereich.

\textbf{Maschinenepsilon.} Mit $2$ Mantissenbits (normalisiert $1{.}b_1b_2$) ist die
nächste darstellbare Zahl über $1{.}00$ gerade $1{.}01_2 = 1+2^{-2}$, also
\[
  \varepsilon_{\mathrm{mach}} = 2^{-2} = \erg{\tfrac14 = 0{,}25}.
\]

\textbf{Auswertung von $\dfrac{1}{(1+0{,}1)-1}$.} Darstellbare Mantissen in $[1,2)$:
$1{,}00\!=\!1$, $1{,}01\!=\!1{,}25$, $1{,}10\!=\!1{,}5$, $1{,}11\!=\!1{,}75$ (Abstand $0{,}25$).
Runden von $1{,}1$:
\[
  |1{,}1-1{,}0| = 0{,}1 \;<\; |1{,}1-1{,}25| = 0{,}15
  \;\Rightarrow\; \mathrm{fl}(1{,}1)=1{,}0 .
\]
Damit ist der \textbf{Nenner nach Rundung}
\[
  \mathrm{fl}(1{,}1)-1 = 1{,}0-1{,}0 = \erg{0},
  \qquad\text{Endergebnis } \frac{1}{0} = \erg{\infty\ (\text{Overflow})}.
\]
\textbf{Vergleich \& Auslöschung.} Mathematisch ist
$\frac{1}{1{,}1-1}=\frac{1}{0{,}1}=10$. Die beiden fast gleichen Zahlen $1{,}1$ und $1$
werden subtrahiert; da $1{,}1$ schon auf $1{,}0$ gerundet wurde, löschen sich bei der
Subtraktion \emph{alle} signifikanten Stellen aus -- der Nenner wird exakt $0$ statt
$0{,}1$. Der reine Rundungsfehler dominiert das Ergebnis, die anschließende Division
verstärkt ihn katastrophal ($\infty$ statt $10$). Das ist \emph{catastrophic cancellation}.

\teil{b}{Gitter-Suche für das lineare System \hfill 4 P}

System: $\;w+2b=4,\quad 3w+b=5$.

\textbf{Fehlermaß.} Summe der absoluten Residuen (L1) -- misst direkt die Verletzung
\emph{beider} Gleichungen, ist ohne Taschenrechner auswertbar und robust:
\[
  E(w,b) = |\,w+2b-4\,| + |\,3w+b-5\,| .
\]
\textbf{Auswertung über $w,b\in\{0,1,2\}$:}
\begin{center}\small
\begin{tabular}{c|ccc}
\toprule
$E(w,b)$ & $b=0$ & $b=1$ & $b=2$\\
\midrule
$w=0$ & $9$ & $6$ & $3$\\
$w=1$ & $5$ & $2$ & $\mathbf{1}$\\
$w=2$ & $3$ & $2$ & $5$\\
\bottomrule
\end{tabular}
\end{center}
Kleinster Fehler $E=1$ bei \erg{(w^\ast,b^\ast)=(1,2)}.
\emph{(Beispiel $(1,2)$: $|1+4-4|+|3+2-5|=1+0=1$.)}

\textbf{Vergleich mit der exakten Lösung} $w^\star=\tfrac65=1{,}2$, $b^\star=\tfrac75=1{,}4$:
Die Gitter-Suche liefert $(1,2)$ -- $w$ passt grob, $b=2$ liegt aber deutlich neben
$1{,}4$. Die Gitter-Suche \textbf{löst das System nicht exakt}, da die exakte Lösung
$(1{,}2;\,1{,}4)$ \emph{nicht} auf dem Gitter $\{0,1,2\}^2$ liegt.

\textbf{Verbesserung.} Feineres Gitter / kleinere Schrittweite, am besten adaptiv um den
besten Gitterpunkt verdichten. Eine \emph{exakte} Lösung findet die Methode nur, wenn die
Lösung zufällig auf dem Gitter liegt -- mit endlichem Gitter im Allgemeinen also nicht,
aber durch Verfeinerung beliebig genau approximierbar (exakt löst man besser direkt, z.\,B.\
per Elimination).

\teil{c}{Fehleranalyse von $g(x)=(x-\tfrac32)^2$ auf $[0,4]$ \hfill 8 P}

Gitter $h_1=1$: $\{0,1,2,3,4\}$; \quad Gitter $h_2=\tfrac12$: $\{0;0{,}5;1;1{,}5;\dots;4\}$.
$\hat g$ wird am nächsten Gitterpunkt ausgewertet; Störung $\tilde x = x\pm\tfrac12$.

\textbf{Kondition} $\kappa(x)=|g(x)-g(\tilde x)|$ (Eigenschaft des \emph{Problems}, vom
Gitter unabhängig):
\[
  \kappa(1{,}5)=|g(1{,}5)-g(2)| = |0-0{,}25| = \erg{0{,}25}\ (\text{gut, flach}),
\]
\[
  \kappa(4)=|g(4)-g(3{,}5)| = |6{,}25-4| = \erg{2{,}25}\ (\text{schlecht, steil}).
\]
Im flachen Bereich um das Minimum $x=1{,}5$ ist $\kappa$ klein (gut konditioniert), im
steilen Randbereich groß (schlecht konditioniert) -- $\kappa$ folgt der Steigung von $g$.

\textbf{Stabilität} der diskretisierten $\hat g$: $|\hat g(\tilde x)-\hat g(x)|\le s\,|\tilde x-x|$,
$|\tilde x-x|=0{,}5$. Referenzpunkt $x=2$:
\begin{itemize}\setlength{\itemsep}{1pt}
  \item $h_1=1$: Die Zelle um $x=2$ ist $(1{,}5;\,2{,}5)$, dort $\hat g\equiv g(2)=0{,}25$.
        Die Störung $\tilde x=2\pm0{,}5$ erreicht gerade die Zellgrenzen; im Zellinneren
        ändert sich nichts $\Rightarrow s\approx 0$, sehr stabil.
  \item $h_2=\tfrac12$: Zelle um $x=2$ ist $(1{,}75;\,2{,}25)$. $\tilde x=2{,}5\to \hat g=g(2{,}5)=1{,}0$
        (Sprung $0{,}75$), $\tilde x=1{,}5\to \hat g=g(1{,}5)=0$ (Sprung $0{,}25$).
        $s=\tfrac{0{,}75}{0{,}5}=\erg{1{,}5}$, weniger stabil.
\end{itemize}
$\Rightarrow$ \textbf{$h_1=1$ ist stabiler}: das grobe Gitter (Zellbreite $1>0{,}5$) fängt
die Störung innerhalb einer Zelle ab; beim feinen Gitter (Zellbreite $0{,}5$) überspringt
die Störung Zellen und erzeugt Sprünge.

\textbf{Konsistenz} $|\hat g(x)-g(x)|\le c$ -- am Minimum:
\begin{itemize}\setlength{\itemsep}{1pt}
  \item $h_1=1$: $g$-Werte $\{2{,}25;\,0{,}25;\,0{,}25;\,2{,}25;\,6{,}25\}$, diskretes
        Minimum $0{,}25$ (bei $x=1,2$). Wahres Minimum $g(1{,}5)=0$ liegt \emph{zwischen}
        den Gitterpunkten $\Rightarrow c=|0{,}25-0|=\erg{0{,}25}$.
  \item $h_2=\tfrac12$: $1{,}5$ ist Gitterpunkt, $\hat g(1{,}5)=g(1{,}5)=0$ $\Rightarrow c=\erg{0}$.
\end{itemize}
$\Rightarrow$ \textbf{$h_2$ ist konsistenter}: das feinere Gitter trifft das Minimum exakt.

\textbf{Konvergenz} (Stabilität $+$ Konsistenz nahe dem Minimum):
$h_1$ hat Konsistenzfehler $0{,}25$ (verfehlt das Minimum), ist aber stabil gegen die
Störung; $h_2$ trifft das Minimum exakt (Konsistenzfehler $0$), ist aber instabiler
(unter $\pm0{,}5$ springt $\hat g$ um bis zu $0{,}75$). \textbf{Eine kleinere Schrittweite
führt also nicht zwingend zu einer besseren Approximation}: der Konsistenzgewinn wird bei
verrauschter Eingabe durch den Stabilitätsverlust teils aufgehoben. Konvergenz verlangt
die Balance beider -- der klassische Konsistenz-Stabilität-Trade-off.

% =====================================================================
\aufg{2}{Newton-Verfahren \& raue Verlustlandschaft \hfill 9 Punkte}

\teil{a}{Newton auf $f(x)=x^3-3x+3$, $x_0=\tfrac32$ \hfill 7 P}

Ableitung $f'(x)=3x^2-3$.

\textbf{Iteration 1.}
\[
  f(\tfrac32)=\tfrac{27}{8}-\tfrac92+3=\tfrac{27-36+24}{8}=\tfrac{15}{8},\qquad
  f'(\tfrac32)=3\cdot\tfrac94-3=\tfrac{27-12}{4}=\tfrac{15}{4}.
\]
\[
  x_1=\tfrac32-\frac{15/8}{15/4}=\tfrac32-\tfrac12=\erg{1}.
\]
\textbf{Iteration 2.}
\[
  f(1)=1-3+3=1,\qquad f'(1)=3-3=0
  \;\Rightarrow\; x_2 = 1-\frac{1}{0} = \erg{\text{undefiniert}}.
\]
\textbf{Beobachtung.} Newton landet \emph{exakt} im kritischen Punkt $x_1=1$ mit $f'(1)=0$
(waagrechte Tangente). Die Tangente schneidet die $x$-Achse nie, der nächste Schritt
divergiert -- \textbf{Versagensmodus: verschwindende Ableitung}. (Die einzige reelle
Nullstelle liegt bei $x\approx-2{,}10$ und wird so nie erreicht.)

\textbf{Anderer Punkt mit ähnlicher Beobachtung.} $f'(x)=3x^2-3=0\Rightarrow x=\pm1$; also
\erg{x=-1} (das zweite Extremum, ein lokales Maximum) -- auch dort $f'=0$.

\textbf{Konvergenzkriterium.} $|x_{n+1}-x_n|<\varepsilon$ (bzw.\ $|f(x_n)|<\varepsilon$);
Newton konvergiert (quadratisch), wenn der Start hinreichend nahe an einer \emph{einfachen}
Nullstelle mit $f'\neq0$ liegt.

\textbf{Numerische Approximation von $f'(x_0)$.} Vorwärtsdifferenz
\[
  D_h(x_0)=\frac{f(x_0+h)-f(x_0)}{h}.
\]
Aus der Taylor-Entwicklung $f(x_0+h)=f(x_0)+f'(x_0)h+\tfrac{h^2}{2}f''(x_0)+\dots$ folgt
\[
  D_h(x_0)=f'(x_0)+\tfrac{h}{2}f''(x_0)+\mathcal O(h^2),
  \quad\text{Approximationsfehler } \erg{\approx \tfrac{h}{2}\,f''(x_0)=\mathcal O(h)}
\]
(eine \emph{zentrale} Differenz wäre $\mathcal O(h^2)$).

\textbf{Wie ändert sich die Beobachtung?} Die Differenzenquotient-Steigung mittelt über
$[x,x+h]$ und ist auch nahe $x=1$ in der Regel $\neq0$. Damit entfällt die exakte Division
durch $0$ -- das Verfahren macht einen endlichen Schritt statt zu divergieren. Daraus
abgeleitet: das \textbf{Sekantenverfahren}
\[
  x_{n+1}=x_n-f(x_n)\,\frac{x_n-x_{n-1}}{f(x_n)-f(x_{n-1})},
\]
ableitungsfrei; es ersetzt $f'$ durch die Sekantensteigung aus zwei Punkten und umgeht so
den Versagensmodus der (analytisch) verschwindenden Ableitung.

\teil{b}{Raue Landschaft $f(x)=(x-1)^2+\tfrac12\sin(2\pi x)$ auf $[0,2]$ \hfill 2 P}

\textbf{Warum problematisch?} Der hochfrequente Sinus-Term überlagert die Parabel mit
vielen kleinen lokalen Minima. Der Gradient wechselt häufig das Vorzeichen; ein
\emph{lokaler} Optimierer (Gradientenabstieg/Newton) bleibt im nächstgelegenen kleinen Tal
hängen und erreicht das globale Minimum (nahe $x=1$) nicht.

\textbf{Weiterer Ansatz \& Schrittweite.} Glättung der Verlustlandschaft durch numerische
Integration/Mittelung (gleitender Mittelwert bzw.\ Faltung über ein Fenster). Geeignete
Schrittweite ist die \emph{Periode} des Stör-Sinus: $\sin(2\pi x)$ hat
$T=\tfrac{2\pi}{2\pi}=1$, also \erg{h=1} (oder ein ganzzahliges Vielfaches), damit sich der
Sinus-Term über das Fenster zu null mittelt und nur die glatte Parabel übrig bleibt.

\end{document}
